%=====================================================================================
\section{What can and cannot be tested}\label{sec:cannot}

The central question, whether Fable's market link comes from reading or from remembering, has a clean
answer only if some statements are unknown to the model. This section goes through the ways to get such
statements and explains why none of them works with the data at hand.

\subsection{Waiting for new statements}

The policy corpus has \R{post.pairs} statement pairs from 2024 to July 2026, \R{post.d2.hold.n} of them at
hold meetings. The surprise series ends in 2023, so only the daily yields can be used. No score shows a
significant relation (Appendix Table~\ref{tab:post}): on hold meetings Fable's $t$-statistics are
\R{post.fable.d2.hold.t} ($d_2$) and \R{post.fable.d10.hold.t} ($d_{10}$), Gemini's
\R{post.gemini.d2.hold.t} and \R{post.gemini.d10.hold.t}.

This is not evidence against the pre-2024 result, because the test has almost no power. If Fable's
pre-2024 $R^2$ were the true effect, the current post-2024 sample would detect it with a probability of
\R{power.min} to \R{power.max} (Appendix Table~\ref{tab:power}). Reaching 80\% power would take between
\R{power.d10.hold.need} and \R{power.d10.all.need} meetings, or \R{power.d2.all.years} to
\R{power.d10.all.years} more years at the current pace. These $R^2$ values come from the sample that
first suggested the effect and are, if anything, too high, so the true requirement is likely larger. Worse, Section~\ref{sec:cutoff}
showed that Fable knows the 2024--2025 statements. At most the four statements from March to July 2026
may be unseen.

\subsection{A model with a documented training cutoff}

GPT-4 is listed with a training cutoff of September 2021. In principle it allows a comparison within one
model: statements before the cutoff it may have memorised, statements after it cannot have. In practice
it fails as an instrument. Its \R{gpt4.nscores} scores take only \R{gpt4.distinct} distinct values;
\R{gpt4.atminus1} of them (\R{gpt4.atminus1pct}) are exactly $-1$; \R{gpt4.zeropct} of its hold-meeting
dScores are exactly zero (Fable: \R{fable.zero} of \R{fable.zeron}); and \R{gpt4.zlbfloor} of the
\R{gpt4.zlbn} zero-lower-bound statements sit at the floor of $-1$. A score that rarely changes between
meetings cannot track meeting-to-meeting surprises.

Its regressions show no positive relation before the cutoff (hold-meeting $t$ of \R{h.gpt4.MPS.hold-pre.t}
for MPS and \R{h.gpt4.d2.hold-pre.t} for $d_2$, the latter with the wrong sign) or after it
($|t|\le\R{h.gpt4.d10.hold-post.t}$ on \R{h.gpt4.d2.all-post.n} or fewer statements; Appendix
Table~\ref{tab:gpt4}). On the same post-cutoff statements Fable's $t$-statistics are between
\R{h.fable.d10.all-post.t} and \R{h.fable.d10.hold-post.t}, but Fable may know those statements. GPT-4's
null result therefore says nothing about whether having seen the statements is enough to produce a
market link. It shows only that GPT-4 with this prompt is too coarse to be a control.

\subsection{Recognised against unrecognised statements}

The planned test was to compare the market link on statements that Fable recognises with the link on
statements it does not. Section~\ref{sec:recognition} showed that the second group is nearly empty
(\R{probe.hold.not} meetings): Fable dates
\R{probe.hold.recog} of the \R{probe.hold.n} probed hold meetings exactly, and blinding and paraphrasing do
not change that. Restricting the regressions to the recognised statements simply repeats the main result.

\subsection{What would settle the question}

Three designs could separate reading from remembering. First, chronologically consistent models
\citep{he2025,drinkall2024} are trained only on text available before each date and so cannot remember
what followed a statement. Scoring the corpus with such models, using the same prompts, would give an
uncontaminated version of the market test. Whether such models produce scores fine enough to detect an
effect of this size is an open question; GPT-4's coarse output shows that it cannot be taken for granted.
Second, a second frontier model with a different training cutoff would show whether the link appears only
for statements a model knows. A dating probe on such a model is cheap (a few dollars) but was not run here.
Third, human labels of meeting-to-meeting \emph{changes} would provide a validity anchor that neither
markets nor model memory can touch. A sheet of \R{humansheet.n} randomly drawn hold-meeting pairs is
prepared in the repository but has not been labelled.
